\begin{corollary}[6.2.2]
\label{IV.2.s6.2.2-fr}
Sean $f:X\to Y$ un morfismo plano de preesquemas y $\mathscr E$ un
$\mathcal O_Y$-módulo cuasicoherente. Si
$\dim.\operatorname{proj}(\mathscr E)\leq n$ \textup{(\textbf{0}, 17.2.14)}, se tiene
$\dim.\operatorname{proj}(f^*(\mathscr E))\leq n$.
\end{corollary}

\oldpage[IV-2]{138}

\begin{proposition}[6.2.3]
\EGAReaderTarget{ega4.c04.prop.6-2-3}{IV.2.s6.2.3-fr}
Sean $A$, $B$ dos anillos locales noetherianos, $k$ el cuerpo residual de
$A$, $\varphi:A\to B$ un homomorfismo local y $N$ un $B$-módulo de tipo finito.
Se supone que $B$ y $N$ son $A$-módulos planos. Entonces se tiene
\begin{equation}
\dim.\operatorname{proj}_B(N)
=\dim.\operatorname{proj}_{B\otimes_Ak}(N\otimes_Ak).
\tag{6.2.3.1}\label{IV.2.s6.2.3.1-fr}
\end{equation}
\end{proposition}

Consideremos, en efecto, una resolución izquierda de $N$ por $B$-módulos
libres de tipo finito
\[
\ldots\to L_i\to L_{i-1}\to\ldots\to L_0\to N\to0
\]
Como los $L_i$ y $N$ son $A$-módulos \emph{planos} (siendo $B$ un $A$-módulo
plano), resulta de (2.1.10) que los $Z_i(L_\bullet)$ son $A$-módulos
\emph{planos}, que $L_\bullet\otimes_Ak$ es una resolución izquierda del
$(B\otimes_Ak)$-módulo $N\otimes_Ak$ y que se tiene
$Z_i(L_\bullet\otimes_Ak)=Z_i(L_\bullet)\otimes_Ak$ para todo $i$. Observemos
que, puesto que $B$ es noetheriano, los $Z_i(L_\bullet)$ son $B$-módulos de
tipo finito, de modo que los $Z_i(L_\bullet\otimes_Ak)$ son
$(B\otimes_Ak)$-módulos de tipo finito. Ahora bien, decir que un $B$-módulo
(resp. un $(B\otimes_Ak)$-módulo) de tipo finito es plano equivale a decir que
es \emph{libre} o, también, \emph{proyectivo} $(\mathbf{0}_{\mathrm{III}}, 10.1.3)$. Como los
$Z_i(L_\bullet)$ son $A$-módulos planos, resulta, por otra parte, de
$(\mathbf{0}_{\mathrm{III}}, 10.2.5)$ (donde se hace $C=B$) que, para que
$Z_i(L_\bullet)$ sea un $B$-módulo plano, es necesario y suficiente que
$Z_i(L_\bullet\otimes_Ak)=Z_i(L_\bullet)\otimes_Ak$ sea un $(B\otimes_Ak)$-módulo plano. El menor
entero $n$ tal que $Z_{n-1}(L_\bullet)$ sea un $B$-módulo libre es, por tanto,
también el menor entero tal que $Z_{n-1}(L_\bullet\otimes_Ak)$ sea un
$(B\otimes_Ak)$-módulo libre, lo que demuestra la proposición (\textbf 0, 17.2.1).
